\begin{abstract}
Large-scale unsupervised language models (LMs) are effective at acquiring extensive world knowledge and certain reasoning capabilities, but exerting fine-grained control over their outputs remains challenging because their training lacks supervision. Common techniques to address this challenge involve gathering human judgments comparing different model outputs and then adjusting the LM through fine-tuning so that it better reflects these annotated preferences, with reinforcement learning from human feedback (RLHF) being a widely used approach. Nevertheless, RLHF entails a complicated and sometimes unstable two-stage process: initially fitting a reward model based on human preferences, and subsequently applying reinforcement learning to update the unsupervised LM such that it maximizes this learned reward while staying close to the original parameters. In this study, we propose an alternative reward model parameterization for RLHF that admits a closed-form solution for the optimal policy, thus reframing the conventional RLHF optimization into a straightforward classification objective. The method we introduce, termed \textit{Direct Preference Optimization} (DPO), offers a stable, efficient, and resource-friendly approach by forgoing the need for sampling from the LM during fine-tuning or extensive hyperparameter searches. Our empirical results indicate that DPO is capable of adapting LMs to human preferences on par with or exceeding other prevalent techniques. Specifically, DPO-based fine-tuning surpasses PPO-driven RLHF in sentiment control of generated text, and either matches or improves performance in summarization and single-turn conversational tasks, all while maintaining much greater simplicity in implementation and training procedures.
\end{abstract}

\section{Introduction}
Large unsupervised language models (LMs), when trained on vast corpora, develop a range of remarkable capabilities~\citep{chowdhery2022palm, brown2020language, touvron2023llama,bubeck2023sparks}. Nevertheless, because these models are built using data authored by individuals with differing intentions, expertise, and objectives, not all patterns within the data are desirable to replicate. For instance, although we may require an AI-based code assistant to \textit{recognize} common programming errors so that it can address them, our aim during code generation is for the assistant to emulate the superior coding practices, even if such examples are infrequent in the training corpus. By the same token, we want the language model to \textit{know} about a misconception endorsed by half of the population, yet avoid restating this falsehood as fact in half the instances it is queried about the topic. In essence, it is vital to distinguish and encourage the model's \emph{preferred outputs and behavior} over its full spectrum of \textit{acquired knowledge and capabilities}, thereby enhancing AI safety, effectiveness, and controllability \citep{ouyang2022training}. Existing approaches frequently direct LMs toward reflecting human preferences via reinforcement learning (RL), but as will be shown, the RL objective traditionally adopted can, in fact, be realized precisely through a simple binary cross-entropy objective—this insight considerably streamlines the preference learning framework.

Broadly, prevalent approaches imbue language models with target behaviors by employing curated collections of human preference judgments that typify responses regarded as beneficial or secure. This preference adaptation phase follows an initial unsupervised pre-training regime over large text datasets. While a direct strategy is to apply supervised fine-tuning on selected high-quality human demonstrations, reinforcement learning from human or AI feedback (RLHF/RLAIF; \citep{christiano2017deep,bai2022constitutional}) has established itself as the most effective set of techniques. RLHF relies on constructing a reward model using annotated preference data and subsequently utilizes RL to fine-tune the language model's policy—promoting outputs with high reward values while curbing divergence from the pre-trained model. Despite its success in producing models with sophisticated conversational and programming abilities, the RLHF methodology is substantially more intricate than supervised learning, as it necessitates training several LMs and interactively sampling from the model during optimization, which leads to greater computational demands.

In this study, we present a method for optimizing language models according to human preferences without the need for explicit reward modeling or reinforcement learning frameworks. We introduce \textit{{Direct Preference Optimization} (DPO)}, a novel algorithm that, while implicitly targeting the same optimization objective as conventional RLHF methods—namely, reward maximization under a KL-divergence regularization—is both easier to implement and train. Conceptually, DPO operates by increasing the log probability ratio in favor of preferred over non-preferred outputs. Crucially, it incorporates a dynamic, example-specific importance weighting mechanism that mitigates the risk of model collapse seen with more naive probability ratio formulations. Similar to previous techniques, DPO employs a theoretical framework for modeling preferences (e.g., the Bradley-Terry model; \cite{bradley1952rankanalysis}) to evaluate how well a reward signal reflects observed preference data. Unlike earlier approaches, which leverage such a model to develop a preference-driven loss for reward model training—followed by policy optimization based on the derived reward—DPO reformulates the preference loss directly as a function of the policy using a variable transformation. This enables DPO to learn a policy directly from a dataset of preference-labeled responses, using a straightforward binary cross entropy loss and yielding an optimal policy with respect to an implicit, data-driven reward signal.

The principal contribution of our work is Direct Preference Optimization (DPO), a streamlined, reinforcement-learning-free approach for aligning language models with preference data. Experimental results demonstrate that DPO performs on par with or better than traditional approaches, such as PPO-based RLHF, across a range of preference-sensitive tasks including sentiment adjustment, summarization, and dialogue, even when applied to models with up to 6B parameters.

Self-supervised language models with larger parameter counts are capable of performing certain tasks in a zero-shot manner \citep{radford2019language} or given only a handful of exemplars via prompt engineering \citep{gpt3,megatron,chowdhery2022palm}. Nevertheless, their results on downstream benchmarks and their alignment with the intentions of users can be notably enhanced by performing fine-tuning on curated datasets that contain instructions paired with human-generated completions \citep{mishra-etal-2022-cross,sanh2022multitask,chung2022scaling,thoppilan2022lamda}. This method, commonly referred to as `instruction-tuning,' facilitates large language models’ capacity to generalize beyond the explicit instruction-tuning corpus, resulting in improved overall utility \citep{chung2022scaling}. Although instruction tuning has achieved significant progress, it is typically less labor-intensive to gather \textit{relative} human evaluations of response quality than to solicit full expert-level demonstrations. Accordingly, more recent research has concentrated on further fine-tuning LLMs using datasets built from human preference data, which has led to advances in tasks like translation \citep{kreutzer-etal-2018-reliability}, summarization \citep{stiennon2022learning,ziegler2020finetuning}, narrative generation \citep{ziegler2020finetuning}, and instruction-following \citep{ouyang2022training,ramamurthy2023is}. These approaches first fit a reward model—typically by optimizing a neural network to predict preferences as described by models such as Bradley-Terry \citep{bradley1952rankanalysis}—and subsequently train the language model to maximize this reward through reinforcement learning, often relying on algorithms such as REINFORCE \citep{williams1992reinforce}, proximal policy optimization (PPO; \cite{schulman2017proximal}), or related methods \citep{ramamurthy2023is}. Closely allied research directions include leveraging LLMs trained via instruction-tuning with human feedback to synthesize new preference datasets focused on desirable traits (e.g., safety or harmlessness), relying on weak supervision like textual rubrics provided by humans to guide LLM-generated annotations \citep{bai2022constitutional}. Such strategies represent the convergence of two central research threads: one concerning reinforcement learning-based language model training for varied tasks~\citep{Ranzato2015SequenceLT,paulus2018a,wu2018learning}, and another focused on generic frameworks for learning from human preferences \citep{christiano2017deep,kupcsik2018learning}. While leveraging relative preference judgments is an attractive proposition, applying reinforcement learning to fine-tune large-scale language models presents significant practical hurdles. This work offers a principled and theoretically sound alternative for optimizing relative preferences without recourse to reinforcement learning.

Beyond language-based applications, research on learning policies from preferences has been conducted in both bandit and reinforcement learning contexts, yielding multiple methodologies. When contextual bandit algorithms use preferences or rankings over actions instead of reward signals, the problem is described as a contextual dueling bandit (CDB; \cite{yue2012karmed,dudik2015contextual}). Since explicit rewards are unavailable in these settings, the theoretical treatment of CDBs often replaces the traditional optimal policy with the concept of a \textit{von Neumann winner}—a policy that achieves at least a 50\% expected win rate against any alternative policy \citep{dudik2015contextual}. Notably, CDB algorithms assume that preference annotations are provided in an online manner, whereas preference learning from human feedback typically uses a static, offline collection of preference-labeled action pairs \citep{yan2022human}. Analogously, in \textit{preference-based RL} (PbRL), policies are trained using binary preferences originating from some hidden 'scoring' function rather than from explicit reward signals \citep{BusaFekete2014,ruiz2023dueling}. While numerous PbRL algorithms have been proposed—some of which can leverage off-policy preference datasets—these commonly involve a two-step process: first modeling the underlying scoring (reward) function, and then optimizing policy with respect to it \citep{jain2013learning,BusaFekete2014,christiano2017deep,sadigh2017active,kupcsik2018learning}. Contrasting these frameworks, we propose a unified policy learning approach that directly optimizes the policy to align with preference feedback.

\section{Preliminaries}\label{section:prelims}

We summarize the RLHF pipeline as described by \citeauthor{ziegler2020finetuning} and in subsequent work (\citep{stiennon2022learning, bai2022training, ouyang2022training}). The standard process comprises three major stages: (1) supervised fine-tuning (SFT), (2) preference collection and reward modeling, and (3) reinforcement learning optimization.

\textbf{SFT}: RLHF pipelines typically commence with fine-tuning a pre-trained language model through supervised learning on curated, high-quality data tailored for downstream applications (such as dialogue or summarization), yielding a supervised fine-tuned model $\pi^\text{SFT}$.

\textbf{Reward Modelling Phase}: During the next stage, this supervised model receives prompts $x$ to generate pairs of candidate outputs $(y_1, y_2)\sim \pi^\text{SFT}(y \mid x)$. Human annotators then evaluate these pairs by indicating a preference for one response, notated $y_w\succ y_l \mid x$, where $y_w$ and $y_l$ identify the favored and less favored outputs within each pair. The assumption is that such preferences derive from some latent, inaccessible reward function $r^*(y, x)$. Various techniques are used for preference modeling; the Bradley-Terry (BT) model \cite{bradley1952rankanalysis} is especially prevalent (although more sophisticated approaches such as Plackett-Luce ranking models \citep{plackett1975analysis, luce2012individual} may be used if rankings of multiple responses are available). Under the BT model, the probability distribution over human preferences $p^*$ can be formulated as:

Given a fixed dataset of preference comparisons, denoted as $\mathcal{D}=\bigl\{x^{(i)}, y_w^{(i)}, y_l^{(i)}\bigr\}_{i=1}^N$ and assumed to be sampled from $p^*$, we may introduce a reward model parameterized by $r_{\phi}(x, y)$ and fit its parameters via the principle of maximum likelihood. When treated as a binary classification task, the negative log-likelihood objective can be written as:

\begin{equation}\label{eq:reward_model}
    \mathcal{L}_R(r_{\phi}, \mathcal{D}) = -\mathbb{E}_{(x, y_w, y_l)\sim \mathcal{D}}\bigl[\log \sigma(r_{\phi}(x, y_w)- r_{\phi}(x, y_l))\bigr]
\end{equation}

with $\sigma$ denoting the logistic sigmoid. In the setting of language models, $r_{\phi}(x, y)$ typically builds on top of the SFT model $\pi^\text{SFT}(y \mid x)$, augmenting it with a linear projection over the final layer to yield a scalar reward value, following the conventions in \cite{ziegler2020finetuning}. To mitigate the variance of the learned reward, a common strategy is to normalize the output so that $\mathbb{E}_{x,y\sim \mathcal{D}}\left[r_\phi(x, y)\right] = 0$ holds across $x$.

\textbf{RL Fine-Tuning Phase}: In reinforcement learning, the reward model obtained in the prior step acts as the feedback provider for further language model adjustment. Adopting methods from previous literature~\citep{jaques2017sequence, jaques2020human}, the optimization target takes the following form:

\begin{equation}\label{eq:RL}
\max_{\pi_{\theta}}  \mathbb{E}_{x\sim \mathcal{D}, y\sim \pi_{\theta}(y \mid x)}\bigl[r_{\phi}(x, y)\bigr] - \beta\mathbb{D}_{\textrm{KL}}\bigl[\pi_{\theta}(y\mid x)\mid \mid \pi_\text{ref}(y\mid x)\bigr],
\end{equation}

where $\beta$ is a tunable parameter regulating divergence from the reference policy $\pi_\text{ref}$—often the initial SFT model $\pi^\text{SFT}$. The policy $\pi_\theta$ itself is also typically initialized to $\pi^\text{SFT}$. This KL-regularization is crucial to ensure that the policy does not move too far away from the regions where the reward model's predictions remain reliable, and to discourage a collapse into only a few high-scoring outputs, thus preserving response diversity. Due to language models generating discrete sequences, the resulting optimization problem is not directly differentiable and is thus addressed via reinforcement learning algorithms. A prevalent approach \citep{ziegler2020finetuning, stiennon2022learning, bai2022training, ouyang2022training} constructs the reward as ${r(x, y) = r_{\phi}(x, y) -\beta (\log \pi_{\theta}(y\mid x) - \log \pi_\text{ref}(y\mid x))}$ and employs Proximal Policy Optimization (PPO) \cite{schulman2017proximal} for maximization.

\section{Direct Preference Optimization}\label{sec:DPO}

Seeking to avoid the complexity of reinforcement learning when scaling policy optimization to large language models, we propose an alternative method that leverages preference data directly. Unlike conventional RLHF workflows that involve first learning a reward model and subsequently maximizing it through RL, our method exploits a reward model parametrization that reveals its corresponding optimal policy in analytic form, thus sidestepping the need for iterative RL procedures. Central to our framework is the realization that an explicit mapping exists between reward functions and their optimal policies; by utilizing this, we convert the conventional loss on reward parameters into a policy-centric loss. This change-of-variables perspective permits policy optimization under models like Bradley-Terry without independently fitting a separate reward function. Effectively, the policy network jointly represents the language model and its induced (implicit) reward landscape.

\textbf{Derivation of the DPO objective.} We begin from the standard reinforcement learning objective, Eq.~\ref{eq:RL}, which utilizes an arbitrary reward function $r$. Building on previous studies~\citep{peters2007reinforcement, peng2019advantage, korbak2022reinforcement, go2023aligning}, it is readily verified that the optimal policy corresponding to the KL-regularized reward maximization (Eq.~\ref{eq:RL}) is:

\begin{equation}\label{eq:op_policy}
    \pi_r(y\mid x) = \frac{1}{Z(x)}\pi_\text{ref}(y\mid x)\exp\left(\frac{1}{\beta}r(x, y)\right),
\end{equation}

where the normalization constant, or partition function, $Z(x) =\sum_{y}\pi_\text{ref}(y\mid x)\exp\left(\frac{1}{\beta}r(x, y)\right)$ ensures proper normalization (see Appendix \ref{app:derivation1} for details). Even when we substitute an estimated reward model $r_{\phi}$ for the true $r^*$, computing $Z(x)$ remains a computationally expensive operation~\citep{korbak2022reinforcement, go2023aligning}, thereby limiting the direct utility of this policy formulation. Nonetheless, Eq.~\ref{eq:op_policy} may be manipulated to solve for the reward in terms of the optimal policy $\pi_r$, the reference distribution $\pi_\text{ref}$, and the partition function $Z(\cdot)$. Specifically, applying the logarithm to both sides of Eq.~\ref{eq:op_policy} and rearranging leads to:

\begin{equation}\label{eq:main_eq}
    r(x,y) =\beta \log \frac{\pi_r(y\mid x)}{\pi_\text{ref}(y\mid x)} + \beta \log Z(x).
\end{equation}

We can implement this alternative parameterization using the true reward function $r^*$ and its associated optimal policy $\pi^*$. Conveniently, for the Bradley-Terry model, only reward differences between two outputs matter, as captured by ${p^*(y_1 \succ y_2 \mid x) = \sigma(r^*(x, y_1) - r^*(x, y_2))}$. Replacing $r^*(x,y)$ with the expression from Eq.~\ref{eq:main_eq} in the preference likelihood Eq.~\ref{eq:bradley-terry}, the partition function $Z(x)$ cancels out, allowing the human preference probability to be represented solely through the optimal policy $\pi^*$ and the reference policy $\pi_\text{ref}$. As a result, the optimal RLHF policy $\pi^*$ satisfies the following relationship under the Bradley-Terry framework:

\begin{equation}\label{eq:objective}
    p^*(y_1\succ y_2 \mid x)=\frac{1}{1 + \exp\left(\beta \log \frac{\pi^*(y_2\mid x)}{\pi_\text{ref}(y_2\mid x)} - \beta \log \frac{\pi^*(y_1\mid x)}{\pi_\text{ref}(y_1\mid x)}\right)}
\end{equation}

Appendix~\ref{app:derivation2} presents the step-by-step derivation. While the above uses the Bradley-Terry preference structure, a similar derivation using the more general Plackett-Luce models~\citep{plackett1975analysis, luce2012individual} is provided in Appendix~\ref{app:plackett_luce_models}.

Having related the human preference probability to the optimal policy, rather than to the reward model directly, we can now define a maximum likelihood estimation procedure for a parameterized policy $\pi_\theta$. By analogy with the reward-based approach (see Eq.~\ref{eq:reward_model}), our policy learning objective is as follows:

\begin{equation}\label{eq:optimum_model}
    \mathcal{L}_\text{DPO}(\pi_{\theta}; \pi_\text{ref}) = -\mathbb{E}_{(x, y_w, y_l)\sim \mathcal{D}}\left[\log \sigma \left(\beta \log \frac{\pi_{\theta}(y_w\mid x)}{\pi_\text{ref}(y_w\mid x)} - \beta \log \frac{\pi_{\theta}(y_l\mid x)}

In this approach, we learn an implicit reward through an alternative parametrization, such that its optimal policy coincides with $\pi_\theta$. Notably, this method can be interpreted as fitting a reparametrized Bradley-Terry model, which imparts desirable theoretical characteristics—such as consistency under certain assumptions regarding the distribution of preference data \cite{bong2022generalized}. We expand upon these theoretical aspects of DPO and their connections to related research in Section~\ref{sec:theory}.

\textbf{How does the DPO update operate?} To elucidate the mechanism underlying DPO, it is instructive to inspect the gradient of the loss $\mathcal{L}_\text{DPO}$. The derivative with respect to the model parameters $\theta$ is given by:

\begin{multline*}\label{eq:gradient}
    \nabla_\theta \mathcal{L}_\text{DPO}(\pi_\theta;\pi_\text{ref}) = \\ -\beta\mathbb{E}_{(x, y_w, y_l) \sim \mathcal{D}} \bigg[\underbrace{\sigma(\hat{r}_\theta(x, y_l) - \hat{r}_\theta (x, y_w))}_\text{greater emphasis when reward estimate is inaccurate}\bigg[\underbrace{\nabla_\theta\log \pi(y_w \mid x)}_\text{promote $y_w$} - \underbrace{\nabla_\theta\log\pi(y_l \mid x)}_\text{diminish $y_l$}\bigg]\bigg],
\end{multline*}

where $\hat{r}_\theta(x, y) = \beta \log \frac{\pi_\theta(y \mid x)}{\pi_\text{ref}(y \mid x)}$ specifies the reward, which is implicitly determined by the language model $\pi_\theta$ relative to the reference model $\pi_\text{ref}$ (see also Section~\ref{sec:theory}). At a high level, the gradient encourages the model to increase the probability of the preferred responses $y_w$ and reduce the probability of the less preferred responses $y_l$. The update assigns higher weight to examples where the implicit reward $\hat{r}_\theta$ erroneously assigns higher value to dispreferred completions, scaled by $\beta$, thereby adjusting for the degree to which the implicit reward misranks completions and reflecting the intensity of the KL regularization. Our experimental findings highlight that this adaptive weighting is essential; omitting the weighting factor, as in a simplistic variant, may lead to significant degradation of the language model’s outputs (see Appendix Table~\ref{tab:unlikelihood_generations}).

\textbf{DPO summary.} The standard procedure for DPO consists of: 1) Drawing completions $y_1, y_2 \sim \pi_\text{ref}(\cdot \mid x)$ for each prompt $x$, using human judgments to annotate these samples and form an offline preference dataset $\mathcal{D} = \{x^{(i)}, y_w^{(i)}, y_l^{(i)}\}_{i=1}^N$, and 2) training the language model $\pi_\theta$ by minimizing $\mathcal{L}_\text{DPO}$, given $\pi_\text{ref}$, $\mathcal{D}$, and a target $\beta$. In real-world applications, preference datasets that are already available are often reused rather than creating new completions and collecting additional human annotations. Since most preference data originate from sampling with $\pi^\text{SFT}$, we set $\pi_\text{ref} = \pi^\text{SFT}$ whenever possible. Alternatively, if $\pi^\text{SFT}$ is unavailable, we initialize $\pi_\text{ref}$ by maximizing the likelihood of the preferred completions $(x, y_w)$: specifically, ${\pi_\text{ref} = \argmax_{\pi}\mathbb{E}_{x, y_w \sim \mathcal{D}}\left[\log \pi(y_w \mid x)\right]}$. This approach reduces the distribution mismatch between the true (but unknown) reference distribution and the $\pi_\text{ref}$ employed in DPO. Additional implementation and hyperparameter details are presented in Appendix~\ref{app:implementation}.

\section{Theoretical Analysis of DPO} Here, we provide further insights into the DPO method, offering theoretical justification and drawing connections between the benefits of DPO and challenges faced by actor-critic approaches commonly utilized in RLHF—such as PPO~\cite{schulman2017proximal}.

\label{sec:theory}

\subsection{Your Language Model Is Secretly a Reward Model} DPO unifies reward learning and policy optimization through a single maximum likelihood framework, effectively eliminating the need to explicitly model rewards and subsequently use RL for policy learning. It is important to note that the objective in Eq. \ref{eq:main_eq} aligns with the Bradley-Terry model where rewards are given by $r^*(x, y) = \beta \log\frac{\pi^*_\theta(y \mid x)}{\pi_\text{ref}(y \mid x)}$, and optimizing our parametrized model $\pi_{\theta}$ becomes equivalent to the reward model training described in Eq. \ref{eq:reward_model} under an appropriate reparameterization. This section develops the theoretical foundation supporting this equivalence, demonstrates that this approach does not reduce the expressiveness of the reward models learned, and proves that the optimal policy can be precisely recovered. We start by introducing an equivalence relation among reward functions.

\begin{definition}
We define two reward functions $r(x, y)$ and $r'(x, y)$ to be equivalent if and only if there exists a function $f$ such that ${r(x, y) - r'(x, y) = f(x)}$.
\end{definition}

It can be readily verified that this forms an equivalence relation, which divides the space of reward functions into distinct equivalence classes. The following lemmas summarize key implications:

\begin{lemma}\label{lemma:same_prefrence}
Within the Plackett-Luce framework, including the Bradley-Terry model as a special case, any pair of reward functions that are equivalent generate identical preference distributions.
\end{lemma}

\begin{lemma}\label{lemma:same_policy}
    If two reward functions belong to the same equivalence class, then they yield the same optimal policy for the constrained reinforcement learning problem.
\end{lemma}

Since the proofs are elementary, we provide them in Appendix \ref{app:lemma1}. The first lemma addresses a classical under-specification challenge inherent to the Plackett-Luce model family \cite{plackett1975analysis}. This under-determination necessitates the imposition of additional constraints for identifiability when deriving maximum likelihood estimates as in Eq. \ref{eq:reward_model} \cite{bong2022generalized}. The second lemma asserts that every reward function within a given equivalence class produces the same optimal policy, implying that in our main objective, it suffices to recover any representative reward function from the correct class. The following theorem, proven in Appendix~\ref{app:thm1}, formalizes this property:

\begin{theorem}\label{thm:main}
    Given mild conditions, all equivalence classes of reward functions that align with the Plackett-Luce (and specifically the Bradley-Terry) models can be written in the form ${r(x, y) = \beta \log \frac{\pi(y\mid x)}{\pi_\text{ref}(y\mid x)}}$, where $\pi(y\mid x)$ is an appropriate model and $\pi_\text{ref}(y\mid x)$ denotes a chosen reference model.
\end{theorem}

\begin{sproof}
    Consider any reward function $r(x, y)$ that determines an associated optimal model $\pi_r(y \mid x)$, as defined in Eq. \ref{eq:op_policy}. We will demonstrate that within the equivalence class of $r$, there exists a representative admitting the stated reparameterization. Define the projection operator $f$ by
\begin{equation}
    f(r; \pi_\text{ref}, \beta)(x, y) = r(x, y) - \beta\log\sum_{y}\pi_\text{ref}(y\mid x)\exp\left(\frac{1}{\beta}r(x, y)\right)
\end{equation}
Here, $f$ adjusts $r$ by subtracting the log of the partition function associated with $\pi_r$, which depends solely on $x$. Thus, $f(r; \pi_\text{ref}, \beta)(x, y)$ remains within the equivalence class of $r(x, y)$. Substituting the expression for $r$ from Eq.~\ref{eq:main_eq} (valid for arbitrary reward functions), we see that $f(r; \pi_\text{ref}, \beta)(x, y) = \beta \log \frac{\pi_r(y\mid x)}{\pi_\text{ref}(y\mid x)}$. Therefore, the projection $f$ generates a member of $r$’s equivalence class in the required reparameterized form, indicating that our reward model remains fully expressive under this parameterization.
\end{sproof}

An alternative interpretation of Theorem~\ref{thm:main} is that it pinpoints, for each equivalence class, the specific reward function chosen by the DPO reparameterization—namely, the reward function that fulfills:

\begin{equation}\label{eq:lag_p}
     \sum_{y}\underbrace{\pi_\text{ref}(y\mid x)\exp\left(\frac{1}{\beta}r(x, y)\right)}_{=\pi(y\mid x)\text{, via Thm.~\ref{thm:main} reparam.}} = 1,
\end{equation}

That is, $\pi(y\mid x)$ defines a proper probability distribution (with non-negative probabilities summing to one). By inspecting Eq.~\ref{eq:op_policy}, it becomes clear that Eq.~\ref{eq:lag_p} expresses the partition function corresponding to the optimal policy derived from $r(x, y)$. A central idea behind the DPO approach is that by enforcing targeted constraints within the generally under-determined Plackett-Luce family of preference models—including, specifically, Bradley-Terry—we can maintain the expressive power for modeling rewards while also rendering the optimal policy in Eq. \ref{eq:op_policy} computationally tractable for any prompt $x$.

\subsection{Instability of Actor-Critic Algorithms}
The framework presented here is also useful for identifying sources of instability in standard actor-critic techniques commonly employed in RLHF workflows, such as PPO. We restrict our attention to the RL fine-tuning procedure as described in Section \ref{section:prelims}, and draw parallels to the control as inference perspective on constrained RL, as discussed in \cite{levine2018reinforcement} and formalized in \ref{eq:RL}. Considering a parameterized policy model $\pi_{\theta}(y\mid x)$, we minimize the divergence $\mathbb{D}_{\text{KL}}[\pi_{\theta}(y|x) \mid \mid \pi^*(y\mid x)]$, where $\pi^*$ denotes the optimal policy associated with the reward $r_{\phi}(y, x)$ from Eq. \ref{eq:optimum_model}. After algebraic simplification, this gives rise to the following optimization criterion:

\begin{equation}\label{eq:AC}
    \max_{\pi_{\theta}}\mathbb{E}_{\pi_{\theta}(y\mid x)}\bigg[\underbrace{r_{\phi}(x, y) -\beta\log\sum_{y}\pi_\text{ref}(y\mid x)\exp\left(\frac{1}{\beta}r_{\phi}(x, y)\right)}_{f(r_{\phi}, \pi_\text{ref}, \beta)} - \underbrace{\beta\log\frac{\pi_{\theta}(y\mid x)}{\pi_\text{ref}(y\mid x)}}_{\text{KL}}\bigg]
\end{equation}

This objective aligns with that optimized in previous research  
\citep{ziegler2020finetuning, stiennon2022learning, bai2022training, ouyang2022training}, employing the DPO-equivalent reward for the reward family $r_{\phi}$. Here, the normalization component within $f(r_{\phi}, \pi_\text{ref}, \beta)$ can be seen as the soft value function associated with the reference policy $\pi_\text{ref}$. Although this normalization does not alter the optimal policy, omitting it can result in policy gradients with substantial variance, thus destabilizing the learning process. Addressing the normalization through a trained value function is possible, yet may introduce its own optimization challenges. A common alternative in earlier works involves reward normalization based on a human completion baseline—effectively a single-sample Monte Carlo approximation of the normalization factor. Notably, the DPO reparameterization delivers a reward formulation that obviates the necessity for such baselines.

\section{Experiments}
This section presents an empirical assessment of DPO’s effectiveness for preference-driven policy learning. We begin by examining a controlled text-generation scenario, posing the question: how does DPO balance reward maximization against KL-divergence minimization with the reference policy, particularly in comparison to standard preference-based algorithms like PPO? Subsequently, we analyze DPO’s capability on larger-scale models and more challenging RLHF benchmarks, such as summarization and dialogue tasks. Our findings indicate that, even with minimal hyperparameter tuning, DPO consistently matches or exceeds the performance of robust approaches—including PPO-based RLHF and choosing the highest-scoring trajectory from $N$ samples using a learned reward model. Prior to reporting empirical results, we outline our experimental design; further details are provided in Appendix~\ref{app:exp_details}.

\textbf{Tasks.} We conduct experiments across three distinct open-ended text generation scenarios. In each case, algorithms are tasked with learning a policy from a preference dataset $\mathcal{D}=\bigl\{x^{(i)}, y_w^{(i)}, y_l^{(i)}\bigr\}_{i=1}^N$. The first scenario, \textbf{controlled sentiment generation}, presents $x$ as the beginning segment of a movie review sourced from the IMDb dataset \cite{maas-EtAl:2011:ACL-HLT2011}, and requires the policy to generate $y$ with explicitly positive sentiment. To ensure an unbiased assessment, we create preference pairs for generations by leveraging a pre-trained sentiment classifier, ensuring that $p(\text{positive}\mid x,y_w)>p(\text{positive}\mid x,y_l)$. For SFT, we further adapt GPT-2-large by fine-tuning on the training portion of the IMDB reviews dataset, proceeding until convergence (additional information provided in App~\ref{app:sentiment_details}). The second task, \textbf{summarization}, uses $x$ as an original Reddit forum post, and expects the policy to produce a summary $y$ distilling the main points. Consistent with existing literature, we use the Reddit TL;DR summarization dataset \citep{volske-etal-2017-tl} in conjunction with human preference annotations collected by \citeauthor{stiennon2022learning}. The SFT model employed here is fine-tuned on summaries authored by humans for forum posts,\footnote{\url{https://huggingface.co/CarperAI/openai_summarize_tldr_sft}} with RLHF implemented via the TRLX framework \citep{leandro_von_werra_2023_7790115}. The preference annotations originate from work by \citeauthor{stiennon2022learning}, who evaluated samples from a different, yet similarly-trained, SFT model. Lastly, the \textbf{single-turn dialogue} setting takes $x$ to be a user-issued query, potentially ranging from scientific inquiries to personal advice requests. The policy must generate a response $y$ that is both engaging and informative; this is evaluated with the Anthropic Helpful and Harmless dialogue dataset \citep{bai2022training}, containing 170k human-assistant exchanges. Each transcript concludes with a pair of model-generated responses, each labeled to indicate human preference. In this context, since no prior SFT model exists, we derive the SFT model by fine-tuning a publicly available language model solely on the favored completions.

\textbf{Evaluation.} We adopt two distinct strategies for evaluating our models. To directly assess each algorithm's capacity to optimize the constrained reward maximization goal, we utilize a controlled sentiment generation scenario. Here, the evaluation is based on each method's Pareto frontier between achieved reward and KL-divergence from the base policy; this analysis is feasible since the true reward function (provided by a sentiment classifier) is available. Conversely, in practical settings where the actual reward function is unknown, we measure performance in terms of \textit{win rate} relative to a baseline, employing GPT-4 as an automated proxy for human judgment. Specifically, in summarization and single-turn dialogue tasks, GPT-4 assesses summary quality and helpfulness of responses. The baseline for summarization consists of the gold standard test set summaries, whereas for dialogue we use the dataset’s preferred response as the benchmark. Prior research suggests language models are more reliable automated evaluators compared to traditional metrics \citep{Chen2023ExploringTU}. To further validate our reliance on GPT-4, we conduct a dedicated human evaluation (Sec.~\ref{sec:human-judgments}), finding a strong correspondence between GPT-4 and human assessments; human-GPT-4 agreement matches or exceeds that between human annotators.

\textbf{Methods.} Beyond DPO, our study includes several established techniques for aligning language model outputs with human preferences. For the summarization domain, we test zero-shot prompts to \textbf{GPT-J} \citep{gpt-j}; for dialogue, we utilize 2-shot prompting with \textbf{Pythia-2.8B} \citep{biderman2023pythia}. We also evaluate an \textbf{SFT} model alongside \textbf{Preferred-FT}, which involves fine-tuning on the selected output $y_w$—sourced from the SFT model for sentiment and summarization, or a generic LM in single-turn dialogue—via supervised learning. As another form of pseudo-supervised learning, we include the \textbf{Unlikelihood} approach~\citep{welleck2019neural}, where the policy is updated to favor $y_w$ while actively reducing the likelihood of $y_l$, modulated by an optional coefficient $\alpha\in[0,1]$ applied to the unlikelihood term. Additionally, we examine \textbf{PPO} \citep{schulman2017proximal} using a reward model trained from preference signals, and \textbf{PPO-GT}, which represents an oracle approach leveraging the exact reward in the sentiment-controlled context. Our experiments with sentiment involve two PPO-GT variants: an off-the-shelf implementation \cite{leandro_von_werra_2023_7790115} and a modified version that incorporates reward normalization and improved hyperparameter tuning for enhanced outcomes; these adaptations are similarly applied to standard PPO with learned rewards. Finally, we test the \textbf{Best of $N$} strategy, which generates $N$ completions per input from the SFT (or Preferred-FT, for dialogue) and selects the top candidate according to a reward model learned from preferences. Despite yielding strong results, this approach is computationally expensive for all but small $N$, as it necessitates $N$ candidate generations per input during evaluation.

\subsection{How well can DPO optimize the RLHF objective?}

The KL-constrained reward maximization objective commonly employed in RLHF methodologies seeks to balance the maximization of reward with a penalty on divergence from a fixed reference policy. Consequently, when evaluating and comparing algorithms, it is essential to consider both the obtained reward and the degree of KL divergence; it is often not ideal to attain marginally improved reward at the cost of substantially larger KL. Figure~\ref{fig:frontier-tldr-main} presents the reward-KL tradeoff for different algorithms under the sentiment task. For each approach, we conduct multiple training runs, varying the policy conservativeness hyperparameter in each run (target KL values $\in\{3,6,9,12\}$ for PPO; $\beta \in \{0.05,0.1,1,5\}$, $\alpha\in\{0.05,0.1,0.5,1\}$ for unlikelihood; and different random seeds for preferred-FT), yielding a total of 22 runs. At intervals of 100 training steps until convergence, we assess the policies using a fixed set of test prompts, reporting both the mean reward measured according to the ground-truth reward function, and the mean sequence-level KL\footnote{Specifically, the sum of KL divergences across all time steps.} to the reference policy $\text{KL}\left(\pi\mid \mid \pi_\text{ref}\right)$. The results indicate that DPO defines a far more efficient reward-KL frontier, delivering superior reward without incurring large KL values. This is remarkable for several reasons: although both DPO and PPO optimize identical objectives, DPO demonstrates significantly better efficiency, as DPO's tradeoff between reward and KL is strictly preferable to PPO’s. Furthermore, DPO surpasses PPO’s performance even when the latter is supplied with the ground truth rewards (PPO-GT).

\subsection{Can DPO scale to real preference datasets?}
\label{sec:dpo-real-datasets}

We proceed to assess the effectiveness of DPO for fine-tuning in tasks involving summarization and single-turn dialogue. Notably, automatic metrics for summarization, such as ROUGE, often do not align well with human preference judgments~\citep{stiennon2022learning}, and previous research has demonstrated that leveraging PPO-based fine-tuning on human preferences leads to more desirable summary outputs. Our evaluation is conducted by generating samples from the test set of the TL;DR summarization dataset and calculating the average rate at which each method's completions are preferred over reference summaries within the test split. For all compared methods, outputs are sampled with temperature values between 0.0 and 1.0, and their respective win rates are depicted in Figure~\ref{fig:frontier-tldr-main} (right). The DPO, PPO, and Preferred-FT methods are each fine-tuned from the identical GPT-J SFT base model\footnote{\url{https://huggingface.co/CarperAI/openai_summarize_tldr_sft}}. Our results indicate that, at a sampling temperature of 0.0, DPO attains an approximate win rate of 61\%, outperforming PPO, which achieves a peak win rate of roughly 57\% at its best temperature (also 0.0). Additionally, DPO secures a higher maximum win rate than the best-of-$N$ selection baseline. It is important to mention that DPO's $\beta$ hyperparameter received minimal tuning, implying the possibility of underestimating DPO’s optimal capacity. Furthermore, DPO demonstrates notably greater resilience to changes in sampling temperature when compared to PPO, whose performance diminishes to that of the original GPT-J model as temperature increases. Preferred-FT yields only marginal gains relative to the SFT baseline. We further provide a direct human evaluation comparison between DPO and PPO in Section~\ref{sec:human-judgments}, observing that DPO outputs at a temperature of 0.25 are chosen over PPO outputs at temperature 0.0 in 58\% of instances.

For single-turn dialogue experiments, we assess the various techniques using a subset of the Anthropic HH test split \citep{bai2022training} that comprises a single round of human-assistant exchange. In GPT-4 based evaluations, we take the test set's preferred completions as the ground truth and calculate each method’s win rate accordingly. Since this benchmark lacks a standardized SFT model, our process begins with the pre-trained Pythia-2.8B; we utilize Preferred-FT to fine-tune a reference model on the designated completions, ensuring output alignment with the desired distribution, and subsequently apply DPO for training. Our analysis also includes a comparison to the Best of 128 Preferred-FT completions—selected since the Best of $N$ approach ceases to improve beyond 128 generations for this particular task (see Appendix Figure~\ref{fig:best-of-n})—as well as to a 2-shot prompted implementation of the base Pythia-2.8B. The results show that DPO achieves comparable or superior performance across optimal temperature settings for each model variant. Additionally, we test an RLHF system, trained via PPO on the Anthropic HH dataset\footnote{\url{https://huggingface.co/reciprocate/ppo_hh_pythia-6B}} from a reputable source\footnote{\url{https://github.com/CarperAI/trlx/tree/main/examples/hh}}, but are unable to identify any prompt or sampling temperature that enables it to outperform the base Pythia-2.8B. Informed by findings on TL;DR and recognizing that both approaches optimize the same reward metric, we treat Best of 128 as a rough estimate of PPO’s upper-bound performance. In summary, DPO stands out as the sole computationally efficient approach yielding improvements over the reference completions on the Anthropic HH dataset, while matching or surpassing the highly resource-intensive Best of 128 strategy. Notably, as Figure~\ref{fig:dialogue-main} illustrates, DPO also reaches its optimal performance in relatively few training steps.

\subsection{Generalization to a new input distribution}

To assess how PPO and DPO perform when exposed to a distribution shift, we apply both policies—trained on Reddit TL;DR summarization—to a different domain: the test set of news articles from the CNN/DailyMail dataset \citep{nallapati-etal-2016-abstractive}. We utilize the optimal sampling temperatures identified in the TL;DR setting (0 and 0.25) for this evaluation. The outcomes are summarized in Table~\ref{tab:ood}. We calculated GPT-4 win rates against the reference summaries provided in the dataset, employing the same GPT-4 (C) prompt used previously for Reddit TL;DR, but with the phrase ``forum post'' changed to ``news article.'' Notably, DPO maintains a substantial advantage over PPO in this new context. These findings offer preliminary support for the hypothesis that DPO policies generalize to new input distributions at least as effectively as PPO policies, despite DPO's lack of access to the extra unlabeled Reddit TL;DR prompts utilized by PPO.

\subsection{Validating GPT-4 judgments with human judgments}
\label{sec:human-judgments}
We undertake a human evaluation to assess the accuracy of GPT-4's preferences, focusing on results from the TL;DR summarization task and employing two distinct prompts for GPT-4. The \textbf{GPT-4 (S)} (simple) prompt asks the model to select the summary that best conveys the main points of the post. In contrast, the \textbf{GPT-4 (C)} (concise) prompt also incorporates a criterion for conciseness, as we observe that with the \textbf{GPT-4 (S)} prompt, GPT-4 tends to favor more verbose and redundant summaries than human raters typically prefer. Complete versions of the prompts are included in Appendix~\ref{app:prompts}. Our comparison spans three conditions, chosen to reflect a spectrum of output qualities: the highest-performing model (DPO, temp. 0.25), the lowest-performing (PPO, temp. 1.0), and an intermediate baseline (SFT, temp. 0.25), all compared against PPO sampled greedily at its optimal temperature. The analysis indicates that GPT-4's decisions, under both prompts, align with human judgment to a degree comparable to inter-human agreement, implying GPT-4 serves as a credible surrogate for human evaluators (multiple human assessments are collected only for the DPO and PPO-1 cases, given limited rater availability). Generally, the \textbf{GPT-4 (C)} prompt yields win rates that more closely mirror human preferences, which is why it is used for the principal outcomes reported in Section~\ref{sec:dpo-real-datasets}. Further methodological information, including the rater interface and the roster of human volunteers, is available in Appendix~\ref{app:human-study}.

\section{Discussion}
Preference-based learning offers an effective and scalable approach for developing capable and aligned language models. In this work, we have presented DPO, a streamlined training strategy that leverages preference data without relying on reinforcement learning techniques. Unlike traditional methods that reshape the preference learning task to fit within an RL framework for compatibility with standard RL algorithms, DPO establishes a correspondence between language model policies and reward mechanisms. This enables the language model to adhere to human preferences \textit{directly} via a straightforward cross-entropy loss function, omitting the need for reinforcement learning and without sacrificing generality. Remarkably, DPO achieves results comparable to or surpassing those of prevalent RLHF methods such as PPO, even with minimal hyperparameter adjustment, thereby substantially lowering the barriers for preference-based language model training.

\textbf{Limitations \& Future Work.} Our findings open several avenues for further investigation. One important direction is understanding how DPO-trained policies generalize to distributions beyond those seen during training, especially when compared to models optimized using explicit reward functions. Preliminary results indicate that DPO achieves similar generalization to PPO-based approaches, but further empirical validation is warranted. Another interesting line of inquiry is whether leveraging self-generated labels from DPO policies facilitates effective utilization of prompts lacking annotations. Additionally, questions remain regarding how over-optimization of the reward signal appears within the direct preference optimization paradigm, and whether the minor performance dip observed in Figure~\ref{fig:dialogue-main}-right exemplifies this phenomenon. While our study involves models up to 6B parameters, extending DPO to much larger, contemporary architectures represents a promising path forward. For model assessment, our analysis reveals that GPT-4-based win rates are sensitive to the choice of prompt; future research could refine techniques for eliciting reliable evaluations from automated agents. Lastly, the principles underlying DPO are not limited to language modeling from human preferences but can potentially benefit the training of generative models across various modalities.